% Part: computability
% Chapter: recursive-functions
% Section: pr-relations

\documentclass[../../../include/open-logic-section]{subfiles}

\begin{document}

\olfileid[fr]{cmp}{rec}{prr}
\olsection{Relations récursives primitives}

\begin{defn}
Une relation $R(\vec x)$ est dite récursive primitive si sa
fonction caractéristique
\[
\Char{R}(\vec x) = \left\{
  \begin{array}{ll}
  1 & \mbox{si $R(\vec x)$} \\
  0 & \mbox{sinon}
  \end{array}
\right.
\]
est récursive primitive.
\end{defn}

Autrement dit, parler d'une relation récursive primitive
$R(\vec x)$ revient à parler de la relation
$\Char{R}(\vec x)=1$, où la fonction récursive primitive
$\Char{R}$ renvoie $1$ ou $0$ pour toute entrée. Par exemple,
la relation $\fn{IsZero}(x)$, qui vaut si et seulement si
$x=0$, correspond à la fonction $\Char{\fn{IsZero}}$
définie par récurrence primitive :
\begin{align*}
\Char{\fn{IsZero}}(0) & = 1,\\
\Char{\fn{IsZero}}(x+1) & = 0.
\end{align*}

On peut aussi composer des relations avec d'autres fonctions
récursives primitives. Les relations suivantes sont donc
récursives primitives :
\begin{enumerate}
\item La relation d'égalité $x=y$, définie par
  $\fn{IsZero}(\left|x-y\right|)$.
\item La relation « inférieur ou égal » $x\leq y$, définie par
  $\fn{IsZero}(x\tsub y)$.\footnote{Correction éditoriale : la source nomme cette relation « less-than », alors que la formule définit bien l'ordre large $x\leq y$.}
\end{enumerate}

\begin{prop}
  L'ensemble des relations récursives primitives est clos
  par les opérations booléennes : si $P(\vec x)$ et
  $Q(\vec x)$ sont récursives primitives, les relations
  suivantes le sont également :
  \begin{enumerate}
  \item $\lnot P(\vec x)$;
  \item $P(\vec x)\land Q(\vec x)$;
  \item $P(\vec x)\lor Q(\vec x)$;
  \item $P(\vec x)\lif Q(\vec x)$.
  \end{enumerate}
\end{prop}

\begin{proof}
  Supposons $P(\vec x)$ et $Q(\vec x)$ récursives
  primitives, c'est-à-dire supposons que leurs fonctions
  caractéristiques $\Char{P}$ et $\Char{Q}$ le soient.
  Il faut montrer que les fonctions caractéristiques
  de $\lnot P(\vec x)$, etc., sont aussi récursives
  primitives. On a
  \[
  \Char{\lnot P}(\vec x) = \begin{cases}
    0 & \text{si $\Char{P}(\vec x)=1$}\\
    1 & \text{sinon}
  \end{cases}
  \]
  et l'on peut définir $\Char{\lnot P}(\vec x)$ par
  $1\tsub\Char{P}(\vec x)$. De même,
  \[
  \Char{P\land Q}(\vec x) = \begin{cases}
    1 & \text{si $\Char{P}(\vec x)=\Char{Q}(\vec x)=1$}\\
    0 & \text{sinon}
  \end{cases}
  \]
  et l'on peut la définir par
  $\Char{P}(\vec x)\cdot\Char{Q}(\vec x)$ ou par
  $\fn{min}(\Char{P}(\vec x),\Char{Q}(\vec x))$.
  Enfin,
  \begin{align*}
    \Char{P \lor Q}(\vec x) & = \fn{max}(\Char{P}(\vec x), \Char{Q}(\vec x)) \text{ et}\\
    \Char{P \lif Q}(\vec x) & = \fn{max}(1 \tsub
  \Char{P}(\vec x), \Char{Q}(\vec x)).
  \end{align*}
\end{proof}

\begin{prop}
  L'ensemble des relations récursives primitives est clos
  par quantification bornée. Si $R(\vec x,z)$ est une
  relation récursive primitive, les relations
  \begin{align*}
    & \bforall{z < y}{R(\vec x, z)} \text{ et}\\
    & \bexists{z < y}{R(\vec x, z)}.
  \end{align*}
  le sont également. La première vaut pour $\vec x$ et
  $y$ si et seulement si $R(\vec x,z)$ vaut pour tout
  $z<y$; la seconde se comprend de même.
\end{prop}

\begin{proof}
  Par convention, $\bforall{z < 0}{R(\vec x,z)}$ est
  vrai, puisqu'il n'existe aucun $z<0$, et
  $\bexists{z < 0}{R(\vec x,z)}$ est faux. Une
  quantification universelle bornée se comporte comme
  un produit fini ou un minimum itéré. Si
  $P(\vec x,y)\defiff\bforall{z < y}{R(\vec x,z)}$,
  on peut donc définir $\Char{P}(\vec x,y)$ par
  \begin{align*}
    \Char{P}(\vec x, 0) & = 1\\
    \Char{P}(\vec x, y+1) & =
    \fn{min}(\Char{P}(\vec x,y),\Char{R}(\vec x,y)).
  \end{align*}
  La quantification existentielle bornée se définit de
  façon analogue au moyen de~$\fn{max}$. On peut aussi
  la déduire de la quantification universelle bornée
  grâce à l'équivalence
  $\bexists{z < y}{R(\vec x,z)}
  \liff\lnot\bforall{z < y}{\lnot R(\vec x,z)}$.
  Notons, par exemple, que
  $\bexists{x \leq y}{\dots x\dots}$
  équivaut à $\bexists{x < y+1}{\dots x\dots}$.
\end{proof}

\begin{prob}
  Montrer que la relation à trois arguments
  $x\equiv y\mod n$ (congruence modulo~$n$)
  est récursive primitive.
\end{prob}

Une autre fonction récursive primitive utile est
la fonction conditionnelle $\fn{cond}(x,y,z)$,
définie par
\begin{align*}
  \fn{cond}(x,y,z) & = \begin{cases}
  y & \text{si $x=0$} \\
  z & \text{sinon}.
\end{cases}
\intertext{Elle se définit récursivement par}
\fn{cond}(0,y,z) & = y,\\
\fn{cond}(x+1,y,z) & = z.
\end{align*}
Pour appliquer le schéma formel de récurrence primitive,
on place l'argument~$x$ en dernière position, puis on
permute les arguments par composition. On peut ainsi
justifier les définitions par cas de fonctions récursives
primitives à partir de relations récursives primitives.

\begin{prop}
Si $g_0(\vec x),\dots,g_m(\vec x)$ sont des fonctions
récursives primitives et
$R_0(\vec x),\dots,R_{m-1}(\vec x)$ des relations
récursives primitives, la fonction $f$ définie par
\[
f(\vec x) = \begin{cases}
    g_0(\vec x) & \text{si $R_0(\vec{x})$} \\
    g_1(\vec x) & \text{si $R_1(\vec{x})$ et non $R_0(\vec{x})$} \\
    \vdots & \\
    g_{m-1}(\vec x) & \text{si $R_{m-1}(\vec{x})$ et si aucune
      des précédentes ne vaut}
    \\
    g_m(\vec x) & \mbox{sinon}
\end{cases}
\]
est également récursive primitive.
\end{prop}

\begin{proof}
  Pour $m=1$, il s'agit simplement de la fonction
  \[
  f(\vec x) = \fn{cond}(\Char{\lnot R_0}(\vec x),g_0(\vec x),g_1(\vec
  x)).
  \]
  Pour $m>1$, il suffit de composer des définitions
  de cette forme.
\end{proof}

\end{document}
